\section{Pattern Matching}
\label{sec:pattern-matching}

Pattern matching finds every instance of a rule's left-hand side in a state, and it is the first step of every expansion. On the Wolfram rule, whose states have at most fourteen edges at depth six, matching, its join and the rewrite together account for $\InstrShareMatch\%$ of executed instructions (Table~\ref{tab:t15-instructions}); canonicalization ($\InstrShareCanon\%$) and the quotient reconstruction ($\InstrShareRecon\%$) are larger.

\subsection{Edge Signatures}

\begin{definition}[Edge Signature]
\label{def:signature}
The \emph{signature} of hyperedge $e = (v_1, \ldots, v_k)$ is the pattern $\sigma(e) = (s_1, \ldots, s_k)$ where $s_i = s_j$ iff $v_i = v_j$.
\end{definition}

For example, edge $(3, 3, 7)$ has signature $(0, 0, 1)$. A pattern edge can match a data edge only
if their signatures are compatible. For a pattern signature with $d$ distinct variables, the number
of compatible data signatures is the Bell number $B_d$~\cite{rota1964number}: $1, 2, 5, 15, 52$ for
$d = 1, \ldots, 5$. The host tests each candidate's signature against the pattern's directly; the
device precomputes the compatible set for each pattern edge (Appendix~\ref{app:gpu-impl}).

\subsection{Candidate Edges by Ancestry}
\label{sec:vertex-index}

A state is a view into the shared edge pool (Section~\ref{sec:storage}), and matching asks two
questions of it: which of its edges contain a given vertex, once a pattern variable is bound, and
which of its edges have a compatible signature, for the first pattern edge. Both are answered from
the state's ancestry. Every edge of a state was added by exactly one state on its chain of
ancestors: the root added all of its edges, and each later state added the edges its event
produced. Each state indexes its own added edges by vertex once, when it is created, so the edges
of $S$ that contain $v$ are
\begin{equation}
    \bigcup_{A \in \text{chain}(S)} \{e \in \text{Contrib}(A)[v] : e \in S\},
\end{equation}
where $\text{Contrib}(A)[v]$ is $A$'s sorted incidence list for $v$, and the test $e \in S$
excludes edges consumed by a later event on the chain. A query costs one search of a small sorted
array per ancestor, nearly every edge it tests is in the state, a state's own index is the size of
its event's right-hand side, and nothing is stored per edge. A global index over all edges of the
evolution, filtered per state, would test every edge ever created at $v$; a state at depth seven
holds under two percent of those.

\subsection{Matching Algorithm}
\label{sec:matching-algorithm}

Algorithm~\ref{alg:matching} gives the task-based matcher.

\begin{algorithm}[htbp]
\caption{Task-Based Pattern Matching}
\label{alg:matching}
\begin{algorithmic}[1]
\Procedure{FindMatches}{$L$, $\hypergraph$}
    \State $e_1, \ldots, e_k \gets \textsc{MatchOrder}(L)$ \Comment{connected: no step is a product}
    \State Enqueue $\textsc{ScanTask}(e_1, \emptyset)$
    \While{task queue non-empty}
        \State $\tau \gets \textsc{Dequeue}()$
        \If{$\tau$ is $\textsc{ScanTask}(e_i, \phi)$}
            \State \textbf{for} each $e \in S$ with $\sigma(e)$ compatible with $\sigma(e_i)$ \textbf{do}
            \State \quad \textbf{if} $\textsc{Compatible}(e, e_i, \phi)$ \textbf{then}
            \State \quad\quad Enqueue $\textsc{ExpandTask}(e_{i+1}, \phi')$
        \ElsIf{$\tau$ is $\textsc{ExpandTask}(e_i, \phi)$ with $i > k$}
            \State \textbf{yield} complete match $\phi$
        \EndIf
    \EndWhile
\EndProcedure
\end{algorithmic}
\end{algorithm}

The match order $e_1, \ldots, e_k$ is fixed when the rule is added, and it is connected: every edge
after the first shares a variable with an earlier one, so no step of the join is a cartesian
product. The host runs it in two forms: a depth-first join that runs a whole search inside one task
and submits a job only for each completed match, and the task form above, which submits a task per
candidate so that work stealing can divide the search. Both forms, and the device, share one
implementation of the recursion, the edge-injectivity check, binding and unbinding, and the choice
of the next pattern edge.

\subsection{Relation to Worst-Case-Optimal Joins}
\label{sec:agm}

Treated as a conjunctive query, a left-hand side has an AGM bound~\cite{atserias2013size}: the
maximum number of results it can have on a database of a given size. A join is worst-case optimal
if its running time matches that bound up to logarithmic factors~\cite{ngo2018worst}. This engine's
edge-at-a-time join meets the bound on left-hand sides of one or two hyperedges, which covers every
rule in the oracle corpus of Section~\ref{sec:benchmarks}; the sweep of Section~\ref{sec:lhs-space}
also covers three and four edges. It does not meet the bound in general: on a cyclic left-hand side
of three hyperedges the bound is $N^{1.5}$ for a state of $N$ edges, while an edge-at-a-time join
performs $\Omega(N^2)$ work, because two of the three edges can be bound before the third
constrains them, whatever the join order. The static analysis of Section~\ref{sec:rule-analysis}
records for each rule whether its left-hand side is acyclic and its integral edge cover, and the
engine warns when a rule's left-hand side is cyclic over three or more edges.

\subsection{Match Forwarding}
\label{sec:forwarding}

When state $s'$ is derived from $s$, a match of $s$ is still a match of $s'$ unless it used an edge
the rewrite removed. The engine forwards such matches from parent to child instead of finding them
again:

\begin{proposition}[Match Forwarding]
Let $\match$ be a match in state $s$, and let $s'$ be derived from $s$ by applying rule $r$ at match $\match'$. Then $\match$ remains valid in $s'$ iff $\text{edges}(\match) \cap \text{edges}(\match') = \emptyset$.
\end{proposition}

Matches reach a child in two ways. A match found in a state is \emph{pushed} to the children the
state already has, and a newly created state \emph{pulls} by walking its ancestor chain and
collecting every ancestor match its own rewrite did not invalidate. Without the push, every state
repeats its ancestors' matching when it is created; without the pull, a match found before a child
existed never reaches it.

A child publishes its parent link before it appears in its parent's child list. The pull treats a
missing parent link as a root and stops. Once a child is visible it can receive a forwarded match,
that match can create a grandchild on another worker, and the grandchild's pull walks up through
the child. If the link were published after the child became visible, that walk could stop early
and miss every match held only by higher ancestors, and nothing would repair it later, because
pulled matches are not stored again in the descendant and each push happens once.

Forwarded matches are collected per parent and delivered to the child together. The static
analysis of Section~\ref{sec:rule-analysis} turns forwarding on for rule sets whose left-hand side
is a join, and a run with a per-state match cap turns it off (Section~\ref{sec:sampling});
Section~\ref{sec:benchmarks} reports the ablation.

\subsection{Delta Matching}
\label{sec:delta}

A match present in the child but not in the parent must use at least one edge the rewrite produced,
because a match built only from surviving edges was already a match in the parent. The child's own
scan therefore looks only for matches that contain a produced edge, and a rewrite produces only the
right-hand side's edges, however large the state is.

Forwarding and the delta scan divide the child's matches between them: forwarding supplies the
matches that use no produced edge, and the delta scan the matches that use at least one. A match
that uses $k$ produced edges is found $k$ times by the delta scan, once from each, so matches are
deduplicated by content, and the duplicate check runs before a match record is allocated.

A validation mode checks this division: it also computes every state's matches by a full scan and
compares them with what forwarding and the delta scan supplied. A missing match is attributed to
forwarding if it uses only surviving edges and to the delta scan otherwise. The validator also
counts how many comparisons it made, so a mismatch count of zero is known to cover a non-zero
number of states.
